%قبل از شروع مسئله ترجیحاً چیزی ننویسید. اما اگر خواستید بنویسید، در قالب لیستی (فقرات و شمارش) بنویسید.
%برای اضافه کردن نمره، بعدی از شروع‌مسئله، [x] را اضافه کنید و به جای x  نمره مورد نظر را قرار دهید. برای زیرمسئله‌ها هم به همین صورت است.
%\شروع‌مسئله[10]

\شروع‌مسئله[9]
یک گراف وزندار و جهتدار با \متن‌لاتین{n} راس و \متن‌لاتین{m} یال داریم. وزن‌ها می‌توانند منفی هم باشند. \متن‌لاتین{s} و \متن‌لاتین{t} دو راس از گراف هستند و \متن‌لاتین{k} عددی طبیعی است. 
\\الگوریتمی از مرتبه زمانی \متن‌لاتین{O(k(n + m))} ارائه کنید که بین تمام مسیرهای از \متن‌لاتین{s} به \متن‌لاتین{t} که دقیقا \متن‌لاتین{k} یال دارند، طول کوتاه‌ترین مسیر را پیدا کند.
\\راهنمایی : یک گراف جدید تشکیل دهید و به ازای هر راس در گراف اصلی مثل \متن‌لاتین{v}، \متن‌لاتین{k+1} راس به شکل \متن‌لاتین{(v, i)} درنظر بگیرید و یال‌ها را طوری قرار دهید که فاصله‌ی \متن‌لاتین{(v, i)} از \متن‌لاتین{(s, 0)} برابر با \متن‌لاتین{i} باشد. 

%اگر مسئله چند قسمت داشت به این صورت عمل کنید:
%\شروع{زیرمسائل}
%\زیرمسئله[0] % مسئلۀ ۱آ
%\زیرمسئله % مسئلۀ ۱ب
%\زیرمسئله % مسئلۀ ۱پ
%\زیرمسئله % مسئلۀ ۱ت
%\پایان{زیرمسائل}

\solution{
از روی گراف اصلی یک گراف جدید می‌سازیم به این شکل که از رئوس این گراف \متن‌لاتین{k + 1} بار کپی می‌گیریم بطوریکه به ازای هر راس مثل \متن‌لاتین{v}، \متن‌لاتین{k + 1} راس در گراف جدید با نام‌های \متن‌لاتین{(v, 0), (v, 1), ..., (v, k)} خواهیم داشت و \متن‌لاتین{k + 1} لایه با شماره‌های ۰ تا \متن‌لاتین{k} از گراف اصلی تشکیل شود. به ازای هریال در گراف اصلی مثل یال \متن‌لاتین{u} به \متن‌لاتین{v}، در گراف جدید \متن‌لاتین{k} یال قرار می‌دهیم که یال \متن‌لاتین{i}ام از راس \متن‌لاتین{(u, i - 1)} به راس \متن‌لاتین{(v, i)}  خواهد بود و وزن آن همان وزن یال \متن‌لاتین{u} به \متن‌لاتین{v} است (یعنی از راس متناظر با \متن‌لاتین{u} در هر لایه به راس متناظر با \متن‌لاتین{v} در لایه‌ی پایین‌تر یال می‌گذاریم). گراف جدید دارای \متن‌لاتین{n(k + 1)} راس و \متن‌لاتین{mk} یال خواهد بود.
\\
 \درج‌تصویر[پهنا=0.3\پهنای‌متن]{F1.PNG}
\\حال ادعا می‌کنیم هر گشت با طول \متن‌لاتین{k} در گراف اصلی از راس \متن‌لاتین{s} به راس \متن‌لاتین{t}، معادل یک مسیر از راس \متن‌لاتین{(s, 0)} در گراف جدید به راس \متن‌لاتین{(t, k)} خواهد بود. به این شکل که از راس \متن‌لاتین{(s, 0)} شروع می‌کنیم و اگر یال \متن‌لاتین{i}ام گشت از \متن‌لاتین{u} به \متن‌لاتین{v}بود، در گراف جدید از راس \متن‌لاتین{(u, i - 1)} به راس \متن‌لاتین{(v, i)} می‌رویم. بنابراین بعد از طی کردن یک گشت به طول \متن‌لاتین{k} در گراف اصلی از \متن‌لاتین{s} به \متن‌لاتین{t}، در گراف جدید یک مسیر از \متن‌لاتین{(s, 0)} به \متن‌لاتین{(t, k)} خواهیم داشت.  برعکس آن هم برقرار است یعنی هر مسیر از \متن‌لاتین{(s, 0)} به \متن‌لاتین{(t, k)} در گراف جدید معادل یک گشت به طول \متن‌لاتین{k} از راس \متن‌لاتین{s} به راس \متن‌لاتین{t} در گراف اصلی خواهد بود به این شکل که اگر یالی مثل \متن‌لاتین{(x, i)} به \متن‌لاتین{(y, i + 1)} را در گراف جدید طی کردیم در گراف اصلی از راس \متن‌لاتین{x} به راس \متن‌لاتین{y} می‌رویم و بنابراین بعد از طی کردن \متن‌لاتین{k} راس \متن‌لاتین{t} می‌رسیم.
\\بنابراین پیدا کردن کوتاه‌ترین گشت با \متن‌لاتین{k} یال از راس \متن‌لاتین{s} به راس \متن‌لاتین{t} در گراف اصلی معادل است با پیدا کردن کوتاه ترین مسیر از راس \متن‌لاتین{(s, 0)} به راس \متن‌لاتین{(t, k)} در گراف اصلی. حال تعریف می‌کنیم: \متن‌لاتین{dis(v, i)} برابر است با کوتاهترین مسیر از راس \متن‌لاتین{(s, 0)} به راس \متن‌لاتین{(v, i)} در گراف جدید. بنابراین
 \متن‌لاتین{dis(v, i) = min(dis(u, i - 1) + w(u, v)}
 به ازای همه‌ی رئوس \متن‌لاتین{(u, i - 1)} که به \متن‌لاتین{(v, i)} یال دارند و \متن‌لاتین{w(u, v)} برابر وزن یال از \متن‌لاتین{u} به \متن‌لاتین{v} است چراکه هر راس از رئوس در لایه‌ی بالاتر \متن‌لاتین{update} می‌شود و فرض می‌کنیم در هرمرحله فاصله رئوس لایه‌های قبلی محاسبه شده‌است و بنابراین فاصله‌ی رئوس این لایه هم به درستی محاسبه می‌شود و برای بدست آوردن مقدار \متن‌لاتین{dis} برای هر راس، به اندازه‌ی تعداد یال‌های ورودی آن راس زمان می‌برد که در مجموع برابر با کل تعداد یال‌ها در گراف جدید می‌شود و از \متن‌لاتین{o(mk + nk)} هزینه می‌کنیم و جواب برابر است با \متن‌لاتین{dis(t, k)}.
\\(برای راه مشابه می‌توان گفت گراف ساخته شده \متن‌لاتین{DAG} است و دور ندارد چراکه یال ها از هر لایه به لایه پایین‌تر است و بنابراین می‌توان یک ترتیب \متن‌لاتین{TOPOLOGICAL} از رئوس با \متن‌لاتین{O(k(n + m))} پیدا کرد که یال‌ها از هر راس به رئوس بعدی است و سپس براساس این ترتیب فاصله هرراس را از روی فاصله‌ی رئوسی که به آن یال دارند بدست می‌آوریم که چون ترتیب \متن‌لاتین{TOPOLOGICAL} است فاصله رئوس همسایه قبلا محاسبه شده اند و با \متن‌لاتین{o(k(n + m))} بدست می‌آید.)
\\
}




%%%%%%%%%%%%%%%%%%%%%%%%%%%%%%%%%%%%%%%%%%%%%%%%%%%%%
% برخی دستورات مفید برای تمرین‌ها. %
%%%%%%%%%%%%%%%%%%%%%%%%%%%%%%%%%%%%%%%%%%%%%%%%%%%%%

% فرمول‌های نمونه:
%$f(x) = \Theta(x)$
%$T(x, y) \leq \log(x) + 2^y + \binom{2n}{n}$
% {\tt code\_function}


% لیست موارد بدون شماره‌گذاری:
%\شروع{فقرات}
%    \فقره مورد اول 
%        \شروع{فقرات}
%            \فقره مورد اول داخلی 
%                \شروع{فقرات}
%                    \فقره اولین مورد خیلی داخلی 
%                    \فقره دومین مورد خیلی داخلی 
%                \پایان{فقرات}
%            \فقره مورد دوم داخلی 
%        \پایان{فقرات}
%    \فقره مورد دوم 
%\پایان{فقرات}

% لیست موارد شماره‌گذاری شده:
%\شروع{شمارش}
%    \فقره یک مورد 
%    \فقره یک مورد دیگر 
%    \فقره باز هم یک مورد دیگر
%\پایان{شمارش}

% برای نوشتن یک فرمول چند خطی که یک شماره بخوره. برای حالت‌های دیگر به سایت زیر مراجعه کنید:
% https://www.overleaf.com/learn/latex/Aligning_equations_with_amsmath
%\begin{align} 
%    \begin{split}
%        (x+y)^3 &= (x+y)^2(x+y) \\
%                &= (x^2+2xy+y^2)(x+y) \\
%                &= (x^3+2x^2y+xy^2) + (x^2y+2xy^2+y^3) \\
%                &= x^3+3x^2y+3xy^2+y^3
%    \end{split}                                 
%\end{align}

% نوشتن ماتریس و آرایه دوبعدی
%\begin{align}
%    A = 
%    \begin{bmatrix}
%        A_{11} & A_{21} \\
%        A_{21} & A_{22}
%    \end{bmatrix}
%\end{align}

% اضافه کردن تصویر با شرح. برای اضافه کردن تصویر بدون شرح، فقط دستور درج‌تصویر را استفاده کنید.
%\شروع{شکل}[h!]
%\تنظیم‌ازوسط
%  \درج‌تصویر[پهنا=0.5\پهنای‌متن]{img.jpg}
%\شرح{شرح تصویر}
%\برچسب{تصویر:۱}
%\پایان{شکل}
%\پایان{زیرمسائل}

%کد زیر را ببینید.
%%\lstinputlisting{source_filename.py}
%\begin{latin}
%\begin{lstlisting}
%# Import the modules
%import sys
%import random
%
%ans = True
%
%while ans:
%question = raw_input("Ask the magic 8 ball a question: (press enter to quit) ")
%
%answers = random.randint(1,8)
%
%if question == "":
%sys.exit()
%
%elif answers == 1:
%print "It is certain"
%
%elif answers == 2:
%print "Outlook good"
%
%elif answers == 3:
%print "You may rely on it"
%
%elif answers == 4:
%print "Ask again later"
%
%elif answers == 5:
%print "Concentrate and ask again"
%
%elif answers == 6:
%print "Reply hazy, try again"
%
%elif answers == 7:
%print "My reply is no"
%
%elif answers == 8:
%print "My sources say no"
%\end{lstlisting}
%\end{latin}